\section{Sensitivity to damping and variance scaling}
\label{sec:mushroom-sensitivity}

We chose the offset $0.01$ in \eqref{eq:conference-damped} as a heuristic in
the controlled pilot. To examine sensitivity to this choice, we declare a
separate comparison of five offsets
$\tau \in \{0.001,0.003,0.01,0.03,0.1\}$ for $s_j = \sqrt{v_j + \tau}$,
together with no scaling, SD ($\tau = 0$), and variance scaling ($s_j = v_j$).
We reuse both sample sizes and all 32 partitions per size, and restrict the
comparison to centered predictors with a fitted intercept. Each fixed rule
receives the same 17 rate factors, two batch sizes, and stopping searches from
Appendix~\ref{sec:mushroom-rate-search}. Comparing the five positive offsets
therefore requires five times the work of evaluating one fixed rule. We do not
select an offset by test loss, and we preserve the earlier protocols and
artifacts.

\begin{table}[ht]
  \centering
  \begin{tabular}{lrr}
    \toprule
    Scaling       &    Aggregate loss &   Rare-group loss \\
    \midrule
    None          & .001341 (.000074) & .005955 (.000888) \\
    SD            & .004476 (.001498) & .004764 (.001079) \\
    Variance      & .311988 (.201772) & .051586 (.008319) \\
    $\tau = .001$ & .001282 (.000054) & .001543 (.000323) \\
    $\tau = .003$ & .001075 (.000031) & .001149 (.000165) \\
    $\tau = .01$  & .000810 (.000039) & .000682 (.000133) \\
    $\tau = .03$  & .000727 (.000044) & .000911 (.000277) \\
    $\tau = .1$   & .000812 (.000053) & .001858 (.000525) \\
    \bottomrule
  \end{tabular}
  \caption{
    Sensitivity to fixed damping offsets: mean held-out half squared loss (MCSE) at
    $n = 2048$, $M = 8192$, under aggregate validation and the original stopping
    checkpoints. Each row receives 34 training paths per split.
  }
  \label{tab:mushroom-sensitivity}
\end{table}

Every tested positive offset reduces rare loss relative to no scaling by more
than two paired MCSE. At offsets $.003$, $.01$, $.03$, and $.1$, both risks
improve by this criterion, including when we add stopping checkpoints. At
$.001$, however, the aggregate difference is $-.000059$ ($.000069$). The
results therefore do not meet the predeclared criterion that both risks should
improve at every positive offset. Offset $.03$ reduces aggregate loss relative
to $.01$ by $.000083$ ($.000019$), while the rare-loss difference is
unresolved. These comparisons neither establish an optimal offset nor cover
untested values.

At $n = 256$, all five positive offsets improve both risks over no scaling at
$M = 8192$ under aggregate validation. Selection by rare validation loss,
however, has a severe failure at $\tau = .03$. Aggregate loss has mean
$42.6229$, median $.004108$, and maximum $1356.8142$, while rare loss has
mean $4.7479$ and maximum $150.2734$. The partition responsible for most of
this loss contains two green training examples and one green validation example.
It selects a finite path with $B = 16$, $M = 2048$, and rate factor $1.4$,
whose rare validation loss is $.000893$. Independent replay passes the declared
numerical tolerance. We retain this failure in every summary.

Variance scaling gives larger losses on the tested grid, but two primary
large-sample selections use the upper rate factor $6.4$. Across caps and
objectives, it reaches the upper boundary in 31 selections, counting repeated
partitions. No other rule reaches a rate boundary in this arm. The search
boundary therefore limits what we can conclude about variance scaling after
tuning. The archive retains 139,264 endpoints, 6,144 selections, 1,501 failed
paths, and 2,403 failed endpoints. We report all paired differences, selected
paths, and results by feature-absence status in
\texttt{notes/mushroom-scaling-sensitivity-results.md}.
